\begin{table*}[ht]
\centering
\small
\begin{tabular}{lccccc}
\toprule
 & VQ-APC & Future-VPC & Masked-NCE & HuBERT & Masked-VPC\\
\midrule
PC/ASR & 8      & 8          & 9           & 8      & 8       \\
f0 & 3      & 3          & 2           & 2      & 3       \\
SV & 4      & 3          & 3           & 3      & 3       \\
\bottomrule
\end{tabular}
\caption{Layers selected for downstream experiments for each model variant.}
\label{tab:layers}
\end{table*}

\section{Appendix}

\subsection{Negative Free Energy of Predictive Coding}
\label{appendix:proof}

The proof of equation \ref{eq:vlb} mirrors those in variational autoencoders \citep{Kingma2014}.
We start with the KL divergence between the posterior distribution $p(z|x_A, x_B)$
and the auxiliary distribution $q(z|x_B)$,
\begin{align}
&\KL\big[q(z|x_B)\|p(z|x_A,x_B)\big] \\
&\quad = \mathbb{E}_{z \sim q} \left[\log \frac{q(z | x_B)}{p(x_B | z) p(z | x_A)} \right] + \log p(x_B|x_A) \notag
\label{eq:proof}
\end{align}
where we assume $x_B \ind x_A \mid z$ or $p(x_B|x_A, z) = p(x_B|z)$.
Because the KL divergence is always positive, we obtain
\begin{align}
&-\log p(x_B | x_A) \\
&\quad \leq \KL \big[q(z|x_B)\|p(z|x_A)\big] - \mathbb{E}_{z \sim q}[\log p(x_B | z)] \notag
\end{align}

\subsection{Future Prediction}

The proof of equation \ref{eq:vlb-ar} involves unrolling $p(x_B, z | x_A)$ or $p(x_{\geq t}, z_{\geq t} | x_{< t})$, where $x_B = x_{\geq t}$ and $x_A = x_{< t}$.
Based on the definition of conditional probability, we have
\begin{align}
&p(x_{\geq t}, z_{\geq t} | x_{< t}) = \prod_{i=t}^T p(x_i, z_i | x_{< i}, z_{t:i-1}) \\
&\quad = \prod_{i=t}^T p(x_i | x_{< i}, z_{t:i}) p(z_i | x_{< i}, z_{t:i-1}) \\
&\quad = \prod_{i=t}^T p(x_i | z_i) p(z_i | x_{< i}),
\end{align}
where the last line makes two reasonable assumptions, $x_i \ind z_{t:i-1} \mid z_i$ and $z_i \ind z_{t:i-1} \mid x_{< i}$.
The variable $z_i$ should represent $x_i$ well without relying on previous latent variables $z_{t:i-1}$.
Given the history $x_{< i}$, computing the representation $z_i$ should not depend on the previous latent variables $z_{t:i-1}$.

%
%

%
%
%
%
%
%
%

%

%
%
%
%
%
%
%
%
%
%
%
%
%

%
%
%
%

\subsection{Pre-Training Recipes}
\label{appendix:recipe}

%
%
%
%
We set a maximum length of 1400 frames per utterance, 
corresponding to about 28 seconds. The learning rate is fixed to $10^{-4}$ under the Adam optimizer,
no learning rate schedule is applied.
We pre-train
\textsc{Base} models on a single A40 with a batch size of 16 for 150 epochs. 
We set model dimension to 768, with inner dimension of feedforward netowrks being 3072.
A dropout probability of 0.1 is applied.

The temperature is set to 1 without annealing for variational training.
In wav2vec 2.0, the similarity value is re-scaled by dividing it with 0.1, 
the temperature for Gumbel-Softmax \citep{jang2016categorical} is annealed from 2 to a minimum of 0.5 by a decay rate of 0.999995,
i.e., the $\tau$ in \Eqref{eq:softmax-q}.
We use 100 negative samples following \citet{baevski2020wav2vec}.

There are a few architectural differences between our setting and \citet{baevski2020wav2vec}.
%
%
We employ a single codebook in our wav2vec 2.0 for quantization.
%
%
%
Rather than Post-LN Transformers, 
we use Pre-LN Transformers for pre-training
and remove the warm-up stage \citep{geiping2023cramming,xiong2020layer}.
We use sinusoidal positional embeddings rather than relative position embeddings used in 
\citet{baevski2020wav2vec}.

%
%
%
%
%
%
%
%
%
%
%
%
%
%
%
%
%
%
%
%
%
%
%
%

\subsection{Downstream Evaluation}
\label{appendix:downstream}

We provide additional details on the experimental setups of downstream tasks. 
The layers chosen for each task is noted in Table~\ref{tab:layers}.

\paragraph{Phone Classification (PC)}
We freeze the pre-trained model and only train a linear layer with a learning rate of $10^{-3}$ for 10 epochs.

\paragraph{Speaker Verification (SV)}

We average frame representations to obtain utterance-level representations, and simply
employ two linear layers to predict 1251 speakers following \citep{Fan2020ExploringW2}. 
The linear classifier is optimized with a learning rate of $10^{-3}$ for 10 epochs.
After training, we take the output of the first linear layer as speaker vectors for speaker verification.
We set the dimension of speaker vectors to 512.

\paragraph{F0 Tracking (F0)}
We train a linear regression layer with a learning rate of $10^{-3}$ for 10 epochs.
We set the minimum and maximum frequency to 50 Hz and 600 Hz respectively.

\paragraph{Automatic Speech Recognition (ASR)}

We use a lightweight sequence-to-sequence encoder for downstream ASR.
The encoder contains two convolutional layers with $(32, 32)$ channels and $(2, 1)$ strides,
and a 4-layer, 256-dimensional bidirectional GRU. The decoder is a unidirectional 256-dimensional GRU.

We adopt a fixed scheduled sampling probability of 0.4 during training. 
We use Adam with learning rates of $10^{-4}$ for all s2s models. 
We employ a dropout rate of 0.2, and a label smoothing rate of 0.1 for regularization. 
%
%
%
We train seq2seq models for 100 epochs, 
and lower the learning rate with a factor of 0.1 for another 20 epochs.

%
%
%
%
%

%
%
%
%
%
%

%
%
%
%
%
%
%
%
%
%

%
%
%
%
%
%
%
%
%
%
%
%
%
%
%
%
%
